\section{Proof of Axis Exactness}
\label{app:axis-exactness}
We prove Lemma~\ref{lem:axis-exactness} for legal read-once components $L,R$
at fixed metadata, under Section~\ref{sec:formal-scope}'s assumptions.
Every logical axis contains an input port, contains at most one port from
each input, and appears at most once in a node's output. These properties
follow because contraction pairs distinct axes from disjoint input subtrees;
permutation reorders axes, and summation removes them from the output.
Removed axes cannot reappear, and separate axes in one child cannot later
merge with each other.

\subsection{Forward Direction}
Suppose $N(L)=N(R)$. The same ordered output classes give the same output
shape. At any output coordinate, both computations select the same input
entries and sum their products over the same remaining indices. Under
exact real arithmetic, the order of these operations does not affect the
result. Thus $L(X)=R(X)$ for every input $X$.

\subsection{Converse Direction}
Suppose $N(L)\ne N(R)$. We construct inputs for which the results differ.

\paragraph{Find a different connection.}
Different output shapes already imply inequivalence. Otherwise, add a port
for each output position and connect it to its class. This makes output
order part of the connection relation. The representations must differ on
some pair of ports: swapping $L,R$ if needed, choose $u,w$ with
$u\sim_L w$ but $u\not\sim_R w$.

\paragraph{Choose inputs.}
In $R$, assign index 1 to the class of $w$ and index 0 to every other class,
including that of $u$. These indices exist because extents are at least two.
Set each input to 1 at the selected coordinate and 0 elsewhere; set scalar
inputs to 1. Call these inputs $X^*$, and let $y$ be the selected output
coordinate.

\paragraph{Compare results.}
In $R$, exactly one summand is 1: changing any summed index selects a zero
input entry, since every class contains an input port. In $L$, a nonzero
summand would require index 0 at $u$ and index 1 at $w$, as fixed by $X^*$
and $y$. Their connection forbids this, so every summand is zero. Hence
\begin{equation}
 L(X^*)[y]=0\ne1=R(X^*)[y].
\end{equation}
Thus equal results for all inputs require equal representations.

\subsection{Example: A Dot Product and Two Independent Sums}
A dot product shares an index; two independent sums do not. With
$A^*=(1,0)$ and $B^*=(0,1)$,
\begin{equation}
\begin{aligned}
 L(A^*,B^*)&=\sum_k A^*_k B^*_k=1\cdot0+0\cdot1=0,\\
 R(A^*,B^*)&=\left(\sum_i A^*_i\right)
             \left(\sum_j B^*_j\right)=(1+0)(0+1)=1.
\end{aligned}
\end{equation}
The extents agree, but the different connections produce different results.
